\documentclass[11pt]{article}

\usepackage[T1]{fontenc}
\usepackage[utf8]{inputenc}
\usepackage{lmodern}
\usepackage[margin=0.86in]{geometry}
\usepackage{microtype}
\usepackage{amsmath,amssymb,amsthm,mathtools}
\usepackage{booktabs,longtable,tabularx}
\usepackage{enumitem}
\usepackage{xcolor}
\usepackage{listings}
\usepackage{xurl}
\usepackage{hyperref}

\definecolor{linkblue}{HTML}{174A7E}
\definecolor{codegray}{HTML}{F3F5F7}
\definecolor{donegreen}{HTML}{176B3A}
\definecolor{opengold}{HTML}{8A5A00}
\definecolor{pendingblue}{HTML}{245B88}
\hypersetup{
  colorlinks=true,
  linkcolor=linkblue,
  citecolor=linkblue,
  urlcolor=linkblue,
  pdftitle={Executable Formalization Blueprint for the McKay Equality},
  pdfauthor={Clawristotle contributors}
}

\lstdefinelanguage{Lean}{
  morekeywords={abbrev,by,class,def,deriving,example,extends,fun,import,
    instance,let,namespace,noncomputable,open,section,structure,theorem,
    variable,where},
  sensitive=true,
  morecomment=[l]{--},
  morecomment=[s]{/-}{-/},
  morestring=[b]"
}
\lstset{
  language=Lean,
  basicstyle=\ttfamily\small,
  backgroundcolor=\color{codegray},
  frame=single,
  framesep=5pt,
  columns=fullflexible,
  keepspaces=true,
  breaklines=true,
  showstringspaces=false
}

\newtheorem{invariant}{Invariant}
\newtheorem{gate}{Compile gate}
\newcommand{\Irr}{\operatorname{Irr}}
\newcommand{\Norm}{\operatorname{N}}
\newcommand{\Cent}{\operatorname{C}}
\newcommand{\Syl}{\operatorname{Syl}}
\newcommand{\IMK}{\mathrm{iMK}}
\newcommand{\done}{\textcolor{donegreen}{\textsc{complete}}}
\newcommand{\openstatus}{\textcolor{opengold}{\textsc{open}}}
\newcommand{\pending}{\textcolor{pendingblue}{\textsc{pending build}}}
\newcommand{\partialstatus}{\textcolor{pendingblue}{\textsc{partial}}}

\title{\textbf{Executable Formalization Blueprint for the McKay Equality}\\
\large Current boundary, proof route, milestones, and compile gates}
\author{Clawristotle contributors}
\date{Audit snapshot: 29 July 2026 (UTC)}

\begin{document}

\maketitle

\begin{abstract}
This is the implementation plan for a kernel-checked Lean proof of the
McKay equality.  It is intentionally not a survey and not a claim that the
proof is finished.  The shortest current route is hybrid: keep the
classification and universal-cover verification required by the published
proof, but replace the projective/Dade-theoretic normal-subgroup seam by the
per-Gallagher-parameter counting equality proved inside Maltempo--Vallejo's
Okuyama--Wajima argument \cite{MaltempoVallejo2026}.  The document records
the exact existing Lean boundary, what is complete, what remains open, and
the compile gate for every major phase.
\end{abstract}

\tableofcontents

\section{Target, trust boundary, and operational invariants}

\subsection{Exact public target}

For a finite group \(G\), a prime \(p\), and \(P\in\Syl_p(G)\), the target is
\[
  |\Irr_{p'}(G)|=|\Irr_{p'}(\Norm_G(P))|.
\]
The proposition already in the repository is
\path{McKayConjecture.Statement}:

\begin{lstlisting}
def Statement (G : Type u) [Finite G] [Group G]
    (p : Nat) [Fact p.Prime] (P : Sylow p G) : Prop :=
  Cardinal.mk (PPrimeIrreducibleCharacter G p) =
    Cardinal.mk
      (PPrimeIrreducibleCharacter (SylowNormalizer P) p)
\end{lstlisting}

The final export should have no mathematical hypothesis beyond this
signature.  The planned endpoint is:

\begin{lstlisting}
theorem McKayConjecture.mckay
    (G : Type u) [Group G] [Finite G]
    (p : Nat) [Fact p.Prime] (P : Sylow p G) :
    McKayConjecture.Statement G p P
\end{lstlisting}

Changing the final theorem name is harmless; weakening its type is not.
In particular, a theorem conditional on a reduction hypothesis, a CFSG
coverage predicate, or family-verification data is an intermediate
milestone and not the deliverable.

\subsection{Pinned environment}

The reproducible environment is:
\[
\begin{array}{ll}
\text{Lean} & \texttt{v4.33.0-rc1},\\
\text{mathlib} &
\texttt{12ab8e82f8447fa639dabe9ffeda74436b72be31}.
\end{array}
\]
The hash is the value in both \path{README.md} and
\path{lake-manifest.json}; no ``latest mathlib'' build is an acceptable
substitute.

\subsection{Trust invariants}

\begin{invariant}[No hidden theorem]
No required reduction, classification, family theorem, or character-table
fact may enter through \texttt{axiom}, \texttt{sorry}, \texttt{admit}, an
empty \texttt{opaque}, a fabricated instance, or inconsistency.  Interface
predicates are permitted only while they remain visibly quantified
hypotheses of intermediate theorems.
\end{invariant}

\begin{invariant}[Reported trusted computing base]
The final audit must run \texttt{\#print axioms} on the exported theorem and
report every dependency.  Standard Lean/mathlib principles such as
propositional extensionality, quotient soundness, and classical choice are
acceptable.  This repository also currently contains 275
\texttt{native\_decide} occurrences in 41 Lean files.  Native decision
procedures therefore form part of the present computational trust boundary;
they must not be described as pure kernel reduction, and generated table
data must still be connected to mathematical objects by checked theorems.
\end{invariant}

\begin{invariant}[Default heartbeats]
Do not set or raise \texttt{maxHeartbeats}.  Split expensive proofs and
large computations into named lemmas and separately compiled files.
\end{invariant}

\begin{invariant}[Exact mathematical boundary]
The numerical shortcut below is used only for the solvable
normal-subgroup seam.  It does not license omitting the quasisimple
inductive-McKay verification or CFSG coverage.
\end{invariant}

\subsection{Build and push invariants}

\begin{enumerate}[label=\textbf{B\arabic*.},leftmargin=*]
  \item A failed \texttt{lake build} is priority zero.  Stop new
  mathematical work, identify the first failing module, restore a full
  build, and only then continue.
  \item Run a targeted module build after each leaf milestone and a full
  \texttt{lake build} after each major phase.  Never start duplicate full
  builds.
  \item A phase is not complete merely because its leaf file builds:
  its umbrella import and the full package must also build.
  \item Before a push, require a green full build, a clean placeholder
  audit, and a reviewed diff.  While work continues and there are changes,
  push on a twelve-hour cadence, never at intervals shorter than twelve
  hours; record the last push time.
  \item LaTeX must build independently of Lean and must not contain claims
  stronger than the declarations it names.
\end{enumerate}

\section{Dated audited status}

\subsection{Snapshot facts}

This audit was taken on 2026-07-29 UTC.  The last commit at audit start was
\texttt{22bbfa2dbbbd044f3cb0883c54d2168bc4204c0f}.  A first unconstrained
dirty-tree \texttt{lake build} exhausted memory while several generated
character-table leaves elaborated concurrently.  Treating that as P0, the
same build was rerun with at most three direct Lean children active and
completed all \(5577\) jobs successfully.  During publication the remote
branch advanced to \texttt{381a0254d78f59b1a85640ba58897eda41fce7ae}
and moved mathlib to the pin recorded above.  After rebasing and fetching
that cache, the \(4010\)-job final-cardinality closure and the complete
\(5583\)-job package both built successfully under the same three-child
memory throttle; the final top-level artifact was written at
2026-07-29 14:24:27 UTC.  Thus P0 is green on the published dependency
state.  The source audit found no live \texttt{sorry}, \texttt{axiom},
\texttt{admit}, or \texttt{maxHeartbeats} declaration in
\path{McKayConjecture/}.

\subsection{Completed/open matrix}

\small
\begin{longtable}{p{0.25\textwidth}p{0.13\textwidth}p{0.54\textwidth}}
\toprule
Layer & Status & Audited fact or remaining obligation\\
\midrule
\endfirsthead
\toprule
Layer & Status & Audited fact or remaining obligation\\
\midrule
\endhead
Public statement and \(p'\)-characters
& \done
& \path{Statement.lean} states the cardinal equality with argument order
\texttt{PPrimeIrreducibleCharacter G p}.\\

Ordinary character, Clifford, Gallagher, and Glauberman APIs used by the
current reduction
& \done
& The relevant implementations and transport lemmas exist and are used by
the central-intersection development.  This row does not claim that every
textbook character-theory theorem has been formalized.\\

Numerical normalizer and central-index induction
& \done
& \path{NumericalReduction.lean} and the central-scalar modules turn local
numerical correspondences into \path{Statement}.\\

Rossi normal-subgroup architecture
& \done
& The central-scalar numerical reduction, join-center step, complement
decomposition, Clifford partitions, and universal-cover/quasisimple
interfaces are implemented.  The stronger exact-normalizer triple
composition still exposes
\path{RestrictedExactNormalizerCompatibility}, but it is off the new
cardinality critical path.\\

Cardinality boundary and adapter
& \done
& The working tree contains the build-verified
\path{CentralScalarCentralIntersectionInertiaFibreCardinalityHypothesis}
and its conversion to the existing fibre-equivalence and final-reduction
interfaces.  It also separates the ordinary count from the automatic
\(q'\)-degree premise and proves their composition into this boundary.\\

Fixed-character and ordinary/\(q'\)-degree bridge
& \partialstatus
& \path{ActionInvariantOver.lean} provides ordinary and \(q'\)-degree
action-invariant lying-over types, explicit finiteness, inner-action
erasure, inertia aliases, and conditional degree-erasure/cardinality
lemmas.  The good-element API and the count itself remain open.\\

Maltempo--Vallejo per-parameter count
& \openstatus
& Equation (2.1), its good-element/Gallagher proof, and its specialization
to the exact Lean inertia fibres have not yet been formalized.\\

Universal-cover and one-Sylow transport
& \done
& Existing APIs reduce a family to one chosen universal cover, one chosen
Sylow subgroup, and only primes dividing the simple target's order.\\

Concrete CFSG coverage
& \openstatus
& \path{FiniteGroupFamily}, \path{CFSGFamilyCoverage}, and
\path{PrimeSupportedCFSGFamilyVerification} are interfaces.  Concrete
finite-Lie-type and sporadic families and an actual classification
coverage proof are absent.\\

\(A_5\)
& \openstatus
& The concrete universal central cover is certified, but the three
exact-normalizer equivariant/projective tables at \(p=2,3,5\) remain
inputs.\\

\(A_6\) ordinary numerical theorem
& \done
& The 31-row ordinary table, exact local normalizer tables, row matchings,
and unconditional numerical McKay statements at \(p=2,3,5\) are proved.\\

\(A_6\) full inductive-McKay datum
& \openstatus
& Three full Sylow-stabilizer equivariance propositions and the rowwise
projective/cohomological compatibility certificates remain; see
Section~\ref{sec:a6}.\\

Alternating groups of degree at least \(7\)
& \openstatus
& The family bookkeeping exists; the prime-supported local
inductive-McKay verification is not complete.\\

Sporadic groups
& \openstatus
& No concrete family verification is present.\\

Lie type outside the final type-\(D\) cases
& \openstatus
& Published theorems are cited in the proof certificate, but no complete
finite-reductive/family implementation is present.\\

Final type \(D\)
& \openstatus
& No \path{FiniteReductive/} or \path{TypeD/} implementation directory
exists.  The local criterion, trichotomy, subgroup construction, extension
maps, and final theorem remain.\\

Universe-polymorphic statement transport
& \done
& The numerical reduction remains specialized to \texttt{Type} (universe
zero), but \path{FinalCardinalityReduction.lean} now exports
universe-polymorphic conditional \path{statement_of_*} wrappers using
\path{Statement.of_universeZero}.  Only their mathematical inputs and the
unconditional public theorem remain open.\\

Unconditional \path{McKayConjecture.mckay}
& \openstatus
& It cannot be exported until the counting boundary, every family
verification, and CFSG coverage are closed and then composed with the
existing universe transport.\\
\bottomrule
\end{longtable}
\normalsize

\section{Proof-route decision}

\subsection{Use the numerical central-scalar reduction}

The efficient target is the already implemented
\path{CentralScalarNumericalInductiveMcKay}, not a new proof of full
inductive McKay for every arbitrary finite group.  The final implemented
spine is intended to be
\[
\begin{gathered}
\text{per-parameter Okuyama--Wajima count}\\
\Downarrow\\
\text{central-intersection inertia-fibre cardinality hypothesis}\\
\Downarrow\\
\text{central-scalar numerical Rossi reduction}\\
\Downarrow\\
\text{quasisimple local data from the CFSG families}\\
\Downarrow\\
\text{\texttt{McKayConjecture.Statement}}.
\end{gathered}
\]
This replaces, at the normal-subgroup seam only, the route through
projective obstruction matching, Dade algebras, and the
Puig--Turull--Ladisch/DGN machinery.  Maltempo--Vallejo explicitly designed
their argument to remove the dependence on Dade's classification in this
solvable-quotient step \cite{MaltempoVallejo2026}.  It does not replace the
family-by-family inductive-McKay verification used by
Cabanes--Sp\"ath \cite{CabanesSpath2026}.

\subsection{Exact specialization of the counting argument}

Fix data quantified by the current cardinality boundary:
\[
\begin{gathered}
X\text{ finite},\quad S\in\Syl_q(X),\quad C\trianglelefteq X,\\
C\cap S\leq Z(X),\qquad
X=C\,\Norm_X(S).
\end{gathered}
\]
Set
\[
\begin{aligned}
R&=O_{q'}(C),\\
L&=RS=CS,\\
H&=\Norm_X(S),\\
R_0&=\Cent_R(S),\qquad \Norm_L(S)=R_0S.
\end{aligned}
\]
In Lean, \(R\) is
\path{ambientNormalIntersectionComplement S C hcentral},
\(L\) is \path{product S C}, \(H\) is
\path{Subgroup.normalizer (S : Set X)}, and the internal copy of
\(\Norm_L(S)=R_0S\) is \path{internalProductNormalizer S C}.
Existing group theory supplies \(R\trianglelefteq X\), the \(q'\)-order of
\(R\), normality of \(L\), \(X=RH\), and the displayed normalizer
description.  Theorem~4.4 coordinates instead use the internal subgroup
\path{productPPrimeKernel S C} of \(L\).  A compile-gated image/equivalence
bridge between that subgroup and the ambient \(R\) is required before any
character identity is stated literally.

Every \(\vartheta\in\Irr_{q'}(L)\) is placed in Gallagher coordinates
\[
  \vartheta=\widehat r\,\mu,
\]
where \(r\in\Irr(R)^S\), \(\widehat r\) is the canonical extension to
\(RS\), and \(\mu\) is linear on \(RS/R\).  In Lean the target parameter
is not literally the same term: it is the image
\[
  \mu^* =
  \path{productLinearQuotientEquiv S C hcentral}\,\mu .
\]
Covariance must show that this transported parameter realizes the paper's
canonical quotient identification.  The existing internal Theorem~4.4
correspondence is intended to be identified literally with
\[
  \vartheta^*=\widehat{r^*}\,\mu^*,
\]
where \(r^*\) is the Glauberman correspondent.

Apply Maltempo--Vallejo with their notation
\[
  A=G=X,\quad K=R,\quad P=S,\quad B=H,\quad C=R_0,
\]
and take their \(Z=1\) and \(\lambda=1\).  Their proof establishes, for
each Gallagher parameter \(\mu\), Equation (2.1):
\[
 \left|\Irr^{\,X}\!\left(X\mid\widehat r\,\mu\right)\right|
 =
 \left|\Irr^{\,H}\!\left(H\mid\widehat{r^*}\,\mu^*\right)\right|.
 \tag{MV-2.1}
\]
Here the superscripts mean invariance under the indicated conjugation
actions.  Since \(X\) acts on itself and \(H\) acts on itself by inner
automorphisms, all ordinary irreducible characters are fixed.

Equation (MV-2.1) counts ordinary irreducible characters, whereas the Lean
endpoint uses \(q'\)-degree subtypes.  This distinction must remain
explicit.  The proof of Theorem~2.4 invokes the standard degree theorem
that characters above the displayed canonical extensions have
\(q'\)-degree \cite[Theorem~5.12]{Navarro2018}.  The formalization must
prove that degree lemma, use a separate ordinary-to-\(q'\) subtype
equivalence, and only then pass to the prime-to-\(q\) subtype.  Ordinary
Clifford transport first identifies Equation (MV-2.1)'s ordinary blocks
with the exact ordinary inertia-fibre types used by Lean.  The checked
degree-erasure adapter then compares those ordinary fibres with their
prime-to-\(q\) subtypes; it does not itself invoke the prime-to-\(q\)
Clifford correspondence.  If that latter correspondence is used earlier
to transport ambient prime-to-\(q\) blocks, its two inertia-index
non-divisibility premises must be supplied explicitly.  Finiteness and
\texttt{Finite.card\_eq.mp} convert the final equality to a nonempty
equivalence when the downstream join-center proof requests one.

\subsection{Why Equation (2.1), not only the statement of Theorem 2.4}

The statement of Maltempo--Vallejo Theorem~2.4 counts all characters above
\(\theta\times\lambda\).  Its proof partitions both sides into disjoint
unions indexed by \(H\)-orbits of linear Gallagher parameters and then
sums the equality over those parameters.  The current Lean boundary is
strictly finer: it asks for the cardinality of each matched Clifford
inertia fibre, for every \(\vartheta\).

An equality of the aggregate disjoint unions does not imply equality of
each summand.  Therefore formalizing only the theorem statement leaves a
logical gap.  The implementation must expose the proof's per-\(\mu\)
Equation (2.1), including its reduction to invariant \(\mu\), and then
use covariance to move from chosen orbit representatives to arbitrary
parameters.  The existing theorem
\path{complementCharacterEquivOfGlauberman_internal_smul} is the expected
target-side covariance bridge.

\subsection{What remains outside the shortcut}

The shortcut applies because the required quotients are \(q\)-solvable;
the automatic-degree and Clifford steps additionally need the exact
\(q'\)-order and index facts listed in P1.E below.  In particular,
\path{productAmbientGlaubermanLeftOuterQuotient_isPPrimeGroup} concerns a
character-inertia quotient and is not a substitute for proving that the
ambient quotient \(X/L\) is \(q'\).  The shortcut does not prove any of
the following:

\begin{itemize}
  \item the inductive-McKay condition for universal covers;
  \item the \(A_6\) projective rows or any other family table;
  \item finite-reductive conditions \(\mathrm A(\infty)\),
  \(\mathrm A(d)\), or \(\mathrm B(d)\);
  \item the final type-\(D\) theorem; or
  \item the CFSG exhaustion.
\end{itemize}

\section{Exact Lean boundary}

The optimized boundary is in namespace
\path{McKayConjecture.InductiveMcKay.NormalSubgroupCentralIntersectionReduction}.
The proof-facing inputs are deliberately separate:

\begin{lstlisting}
def CentralScalarCentralIntersectionOrdinaryInertiaFibreCardinalityHypothesis
    (q : Nat) [Fact q.Prime] : Prop := ...

def CentralScalarCentralIntersectionInertiaFibrePPrimeDegreeHypothesis
    (q : Nat) [Fact q.Prime] : Prop := ...
\end{lstlisting}

Both quantify over \(X,S,C\), the normality, central-intersection, and
generation hypotheses, a
\path{ProductGlaubermanCorrespondence}, and every
\(\vartheta:\)\path{PPrimeIrreducibleCharacter (product S C) q}.  The
first concludes equality of the two ordinary
\path{IrreducibleCharactersOverInertia} cardinalities.  The second says
pointwise that every ordinary character in both fibres has \(q'\)-degree.
The target character is given by
\path{internalTheorem44CharacterEquiv} applied to
\path{complementCharacterEquivOfGlauberman}.

The build-verified adapter chain is:

\begin{lstlisting}
OrdinaryInertiaFibreCardinalityHypothesis
  + InertiaFibrePPrimeDegreeHypothesis
  -> CentralScalarCentralIntersectionInertiaFibreCardinalityHypothesis
  -> CentralScalarCentralIntersectionInertiaFibreEquivHypothesis
  -> CentralScalarCentralIntersectionReductionHypothesis
  -> CentralScalarNormalSubgroupReductionHypothesis
  -> CentralScalarNumericalInductiveMcKay
  -> Statement
\end{lstlisting}

The principal declarations are:

\begin{itemize}
  \item
  \path{centralScalarCentralIntersectionInertiaFibreCardinalityHypothesis_of_ordinary};
  \item
  \path{centralScalarCentralIntersectionInertiaFibreEquivHypothesis_of_cardinality};
  \item
  \path{centralScalarCentralIntersectionReductionHypothesis_of_cardinality};
  \item
  \path{centralScalarNormalSubgroupReductionHypothesis_of_cardinality};
  \item
  \path{centralScalarNumericalInductiveMcKay_of_inertiaFibreCardinality_quasisimple};
  \item
  \path{statement_of_inertiaFibreCardinality_quasisimple};
  \item
  \path{centralScalarNumericalInductiveMcKay_of_inertiaFibreCardinality_universalCover};
  \item
  \path{statement_of_inertiaFibreCardinality_universalCover};
  \item
  \path{statement_of_ordinaryInertiaFibreCardinality_quasisimple};
  and
  \item
  \path{statement_of_ordinaryInertiaFibreCardinality_universalCover}.
\end{itemize}

The classification-side inputs retain their exact names
\path{QuasisimpleInductiveMcKayHypothesis} and
\path{UniversalCoverInductiveMcKayHypothesis}.  They are honest
propositions, not established global theorems.

\section{High-level executable plan}

\begin{longtable}{p{0.08\textwidth}p{0.25\textwidth}p{0.26\textwidth}p{0.31\textwidth}}
\toprule
ID & Phase & Depends on & Phase output\\
\midrule
\endfirsthead
\toprule
ID & Phase & Depends on & Phase output\\
\midrule
\endhead
P0 & Preserve a green baseline & current tree &
Reproducible full build and trust audit.\\
P1 & Close the per-\(\mu\) counting seam & ordinary character,
Clifford, Gallagher, Glauberman APIs &
Closed terms inhabiting the separate ordinary-cardinality and
automatic-degree hypotheses, hence the existing prime-to-\(q\)
cardinality hypothesis.\\
P2 & Finish alternating and finite exceptional local data & exact-normalizer
table interfaces &
Prime-supported family verification, including full \(A_6\) local data.\\
P3 & Formalize finite-reductive and final type \(D\) verification & algebraic
groups, character correspondences, extension maps &
Lie-type one-cover/one-Sylow family verification.\\
P4 & Close CFSG coverage and all family assembly & P2, P3, sporadic
verification &
\path{UniversalCoverInductiveMcKayHypothesis} (and hence the quasisimple
input) without assumptions.\\
P5 & Export and audit the theorem & P1, P4 &
Universe-polymorphic \path{McKayConjecture.mckay} with an audited axiom
report.\\
\bottomrule
\end{longtable}

The dependency order is P0, then P1 and the family work P2--P3 may proceed
in parallel, followed by P4 and P5.  A failure of P0 suspends every other
phase.

\section{P0: baseline, source control, and build discipline}

\subsection*{Low-level plan}

\begin{enumerate}[label=\textbf{P0.\arabic*.},leftmargin=*]
  \item Maintain exactly one current full build.  If it fails, repair the
  first error without starting unrelated proof work.
  \item Record the successful commit/tree hash, Lean toolchain, mathlib
  hash, command, UTC completion time, and axiom/placeholder scan.
  \item Before each leaf edit, identify a smallest target module.  After
  it builds, build its umbrella import.  At a major boundary, run the full
  package build.
  \item Keep generated character-table data in small import leaves.
  Preserve default heartbeats by splitting evaluations rather than raising
  limits.
  \item Push only a reviewed, green state and enforce the twelve-hour
  minimum interval.
\end{enumerate}

\begin{gate}[P0]
\texttt{lake build} exits successfully; the manifest contains the exact
pin above; the source scan has no unfinished declarations or heartbeat
override; no duplicate build is running.
\end{gate}

\section{P1: formalize the per-parameter Okuyama--Wajima count}

\subsection{P1.A: fixed-character and good-element API}

\begin{enumerate}[label=\textbf{P1.A\arabic*.},leftmargin=*]
  \item Define ordinary and \(p'\)-degree action-invariant ``lying over''
  subtypes for an explicit homomorphism into \path{MulAut}.  Keep the
  ordinary subtype primary because Equation (MV-2.1) is an
  ordinary-character statement.
  \item Reuse \path{IrreducibleCharacter.conj_smul} to give ordinary and
  \(p'\)-degree equivalences erasing the invariant field for the
  self-conjugation action.  Add the corresponding inertia aliases and
  \texttt{Nat.card} lemmas in
  \path{McKayConjecture/Character/ActionInvariantOver.lean}.
  \item Give a separate equivalence from the invariant \(p'\)-degree
  subtype to the ordinary invariant subtype under the semantic premise
  that every character in the latter has \(p'\)-degree.  This premise is
  discharged later, not built into Equation (MV-2.1).
  \item Define \(\theta\)-good elements and good quotient conjugacy
  classes.  Prove invariance under conjugacy, replacement inside a normal
  coset, and transport along the character-triple equivalences actually
  used by the proof.
  \item State intermediate lemmas semantically; paper-number aliases may
  be local.  A suitable file is
  \path{McKayConjecture/Character/GoodElement.lean}.
\end{enumerate}

\begin{gate}[P1.A]
The ordinary and \(p'\)-degree fixed-character subtypes are finite; both
inner-action erasure equivalences and the conditional degree-erasure
equivalence build; good-element predicates transport without any
project-specific axiom.
\end{gate}

Items P1.A1--P1.A3 are implemented and targeted-build verified in
\path{ActionInvariantOver.lean}.  Items P1.A4--P1.A5 remain open, so this
gate is only partially closed.

\subsection{P1.B: invariant Gallagher count}

\begin{enumerate}[label=\textbf{P1.B\arabic*.},leftmargin=*]
  \item Formalize the class-function restriction map on representatives of
  good conjugacy classes.
  \item Prove Maltempo--Vallejo Theorem~1.3: fixed characters above an
  invariant normal character are counted by good quotient classes.
  \item Prove the inertia/stabilizer version corresponding to Corollary
  1.4, using the existing Clifford correspondence and its equivariance.
  \item Place these results in
  \path{McKayConjecture/Character/GallagherInvariantCount.lean}.
\end{enumerate}

\begin{gate}[P1.B]
Both invariant and non-invariant normal-character forms compile, and the
latter specializes to the former when the inertia group is top.
\end{gate}

\subsection{P1.C: extension criterion and good-class comparison}

\begin{enumerate}[label=\textbf{P1.C\arabic*.},leftmargin=*]
  \item Formalize the Okuyama--Wajima extension criterion (Theorem~2.1) in
  the exact abelian-quotient strength used here; do not formalize the
  stronger Dade/Ladisch theorem.
  \item Derive Corollary~2.3 by comparing good classes.
  \item Formalize Lemma~2.2, including the \(p'\)-degree version and the
  empty-fibre case.
  \item Put the results in
  \path{McKayConjecture/Character/OkuyamaWajimaInvariantCount.lean}.
\end{enumerate}

\begin{gate}[P1.C]
The equality of fixed-character counts over a Glauberman pair compiles
under exactly the normality, coprimality, and abelian quotient hypotheses
of the published argument.
\end{gate}

\subsection{P1.D1: canonical and covariant extensions}

\begin{enumerate}[label=\textbf{P1.D1.\arabic*.},leftmargin=*]
  \item Use
  \path{CharacterTriple.exists_pPrime_extension_of_normalHall} (and its
  inertia-top form) only for existence.  Its \path{Classical.choose}
  witness is not yet the canonical extension of the paper.
  \item State the determinantal-order characterization of
  \(\widehat r\), prove uniqueness, and then prove covariance under every
  automorphism in the counting action.
  \item Repeat this construction for \(r^*\) and prove compatibility of
  the source and Glauberman-target canonical extensions.
  \item Construct the source quotient parameter \(\mu\).  Define its target
  as
  \path{productLinearQuotientEquiv S C hcentral mu}, and prove that this
  transport is the paper's canonical quotient identification.
\end{enumerate}

\begin{gate}[P1.D1]
The two canonical extensions are uniquely characterized, covariant, and
compatible; the transported target parameter is explicit in every theorem
type.
\end{gate}

\subsection{P1.D2: stabilizers and twisting reductions}

\begin{enumerate}[label=\textbf{P1.D2.\arabic*.},leftmargin=*]
  \item Reduce first to the \(r\)- or \(\theta\)-stabilizer, recording the
  character and action transports supplied by Lemma~2.2.
  \item For a non-invariant \(\mu\), reduce again to its stabilizer and
  identify the two resulting stabilizers with the exact source and target
  character inertia groups.
  \item In the invariant case, twist by the canonical extension and prove
  that the operation commutes with the quotient-parameter transport.
  \item Use covariance to return from chosen stabilizer representatives to
  every \(\mu\), including the empty-fibre cases.
\end{enumerate}

\begin{gate}[P1.D2]
The stabilizer reductions compile as equivalences of the actual
fixed-character fibres; no orbit-representative choice appears in the
exported statement.
\end{gate}

\subsection{P1.D3: Schur--Zassenhaus and good-class comparison}

\begin{enumerate}[label=\textbf{P1.D3.\arabic*.},leftmargin=*]
  \item Choose the Schur--Zassenhaus complement used in Equation (2.2) and
  expose its conjugacy-independence.
  \item Build restriction equivalences to the complement and prove that
  they preserve the relevant good classes.
  \item Apply Corollary~2.3 to obtain Equation (2.2), then transport it back
  across the restriction equivalences.
\end{enumerate}

\begin{gate}[P1.D3]
Equation (2.2) is a semantic equality of fixed-character
\texttt{Nat.card}s and is independent of the chosen complement.
\end{gate}

\subsection{P1.D4: export Equation (2.1)}

\begin{enumerate}[label=\textbf{P1.D4.\arabic*.},leftmargin=*]
  \item Assemble D1--D3 for one explicit source parameter \(\mu\) and its
  transported target parameter \(\mu^*\).
  \item Export Equation (2.1) for every \(\mu\), before taking any
  disjoint-union sum.
  \item Add a regression specialization for an invariant parameter and one
  for a parameter with a proper stabilizer.
\end{enumerate}

\begin{gate}[P1.D4]
A theorem with one explicit Gallagher parameter on each side produces the
ordinary fixed-character cardinality equality.  Its proof does not invoke
the aggregate conclusion of Theorem~2.4 as a converse.
\end{gate}

\subsection{P1.E1: ambient/internal coordinates}

\begin{enumerate}[label=\textbf{P1.E1.\arabic*.},leftmargin=*]
  \item Instantiate \(R,L,H,R_0\) as in the route decision and prove the
  image/equivalence bridge from the ambient
  \path{ambientNormalIntersectionComplement} to the internal
  \path{productPPrimeKernel}.
  \item Reuse
  \path{productPPrimeKernel_isPPrimeGroup},
  \path{productPPrimeKernel_sup_productSylow_eq_top}, and
  \path{productQuotient_isPPrimeGroup}.  The last fact supplies the required
  \(q'\)-quotient \(X/L\); prove the analogous
  \(H/\Norm_L(S)\) fact explicitly.
  \item Put every source \(\vartheta\) in canonical Gallagher coordinates,
  retain its covariance data, and transport \(\mu\) with
  \path{productLinearQuotientEquiv}.
  \item Prove that \path{internalTheorem44CharacterEquiv} is literally the
  character \(\widehat{r^*}\mu^*\) used by Equation (2.1).
\end{enumerate}

\begin{gate}[P1.E1]
The ambient and internal kernels, quotient parameters, canonical
extensions, and the existing Theorem~4.4 target are connected by literal
Lean equalities or named transport equivalences.
\end{gate}

\subsection{P1.E2: automatic degree and inertia indices}

\begin{enumerate}[label=\textbf{P1.E2.\arabic*.},leftmargin=*]
  \item On the source, prove
  \begin{lstlisting}
  not (q ∣ (IrreducibleCharacter.inertia N theta).index)
  \end{lstlisting}
  Derive it from Sylow containment using \path{S.not_dvd_index} and
  \path{Subgroup.index_dvd_of_le}; prove the analogous target statement.
  \item Formalize the exact consequence of Navarro
  Theorem~5.12: above each relevant invariant canonical extension, the
  normal-subgroup degree theorem and the \(q'\)-index force every ordinary
  lying-over character to have \(q'\)-degree.
  \item Apply the theorem on both sides and inhabit
  \path{CentralScalarCentralIntersectionInertiaFibrePPrimeDegreeHypothesis}.
\end{enumerate}

\begin{gate}[P1.E2]
Both pointwise automatic-degree functions compile, as do the two
non-divisibility premises required by the prime-to-\(q\) Clifford API.
\end{gate}

\subsection{P1.E3: Clifford and fixed-action transport}

\begin{enumerate}[label=\textbf{P1.E3.\arabic*.},leftmargin=*]
  \item Specialize Equation (2.1), use inner-action erasure, and identify
  its ordinary blocks with the exact ordinary inertia fibres.
  \item Inhabit
  \path{CentralScalarCentralIntersectionOrdinaryInertiaFibreCardinalityHypothesis}.
  \item Only after P1.E2, use the conditional ordinary-to-\(q'\)
  equivalence on the exact inertia fibres.  This direct adapter is the
  shortest critical path.  If
  \path{CliffordEquivalence.cliffordPPrimeCorrespondence} is used instead
  to transport ambient prime-to-\(q\) blocks, supply its explicit
  inertia-index non-divisibility hypotheses rather than expecting
  typeclass or simplifier inference.
\end{enumerate}

\begin{gate}[P1.E3]
The ordinary Equation (2.1) endpoint and the ordinary Clifford transports
compile for the exact existing source and target characters; the
conditional degree-erasure adapter then reaches the prime-to-\(q\) fibres.
\end{gate}

\subsection{P1.E4: final cardinality assembly}

\begin{enumerate}[label=\textbf{P1.E4.\arabic*.},leftmargin=*]
  \item Apply
  \path{centralScalarCentralIntersectionInertiaFibreCardinalityHypothesis_of_ordinary}
  to the outputs of P1.E2 and P1.E3.
  \item Feed the result through the quasisimple and universal-cover final
  wrappers.  Prefer the direct
  \path{statement_of_ordinaryInertiaFibreCardinality_quasisimple} and
  \path{statement_of_ordinaryInertiaFibreCardinality_universalCover}
  declarations.
  \item Build the targeted integration file, both umbrella imports, and
  the full package.
\end{enumerate}

\begin{gate}[P1.E4]
Closed ordinary-count and automatic-degree terms are accepted by the
direct final wrappers; the targeted file, umbrella import, and full package
all build.
\end{gate}

\section{P2: alternating and finite exceptional families}

\subsection{\(A_6\): exact remaining work}\label{sec:a6}

The completed datum is
\path{alternatingSixAmbientOrdinaryCharacterTableCertificate}.  It gives
the unconditional numerical theorems
\path{alternatingSixCanonicalCover_statement_two_unconditional},
\path{alternatingSixCanonicalCover_statement_three_unconditional}, and
\path{alternatingSixCanonicalCover_statement_five_unconditional}.  These
are not full inductive-McKay data.

\subsubsection*{Low-level plan}

\begin{enumerate}[label=\textbf{P2.A6.\arabic*.},leftmargin=*]
  \item Prove the three propositions (currently definitions, not theorems):
  \path{AlternatingSixTwoCentralBlockRowMatchingEquivariant},
  \path{AlternatingSixThreeDisplayedOrderEquivariant}, and
  \path{AlternatingSixFiveCentralBlockRowMatchingEquivariant}.
  They quantify over the full
  \path{SylowAutStabilizer}, not merely inner automorphisms or the two
  stored outer generators.  Use the proved generator action data plus an
  explicit generation theorem for the stabilizer.
  \item Package those proofs into the existing concrete character-table
  certificates.  At this gate the ordinary equivariant McKay bijections
  are complete, but no projective claim is.
  \item For each matched row and each of \(p=2,3,5\), construct the existing
  propositions
  \path{AlternatingSixTwoConcreteProjectiveCompletion},
  \path{AlternatingSixThreeConcreteProjectiveCompletion}, and
  \path{AlternatingSixFiveConcreteProjectiveCompletion}.
  \item The exact row target is
  \path{ExactNormalizerProjectiveRowData}.  Its witness must provide
  associated projective representations and prove literal
  \path{ProjectiveCompatibilityWitness.factorAgreement} and
  \path{ProjectiveCompatibilityWitness.scalarAgreement}.  If the finite
  calculation first proves equality of \(H^2\)-classes, use the existing
  factor-set cohomology and rescaling API to choose gauges and obtain the
  literal factor equality required by the structure.  Equality of
  cohomology classes alone is not the final field.
  \item Bundle the three completed tables into
  \path{AlternatingSixExceptionalPrimeLocalData}.  Only then is the
  \(A_6\) family verification complete.
\end{enumerate}

\begin{gate}[\(A_6\)]
Construct \path{AlternatingSixExceptionalPrimeLocalData} without
hypotheses.  Verify its three local fields and the derived
prime-supported family theorem by compilation and an axiom report.
\end{gate}

\subsection{\(A_5\), higher alternating groups, and sporadics}

\begin{enumerate}[label=\textbf{P2.F\arabic*.},leftmargin=*]
  \item For \(A_5\), fill the finite exact-normalizer projective tables at
  \(2,3,5\).  The concrete universal central extension is already
  available; do not list the family as complete until the table inputs
  disappear from theorem types.
  \item For alternating degree at least \(7\), formalize the published
  prime-supported equivariant bijections, extension maps, and projective
  compatibility on one cover and one Sylow subgroup per representative.
  \item Define the concrete sporadic family, enumerate its representatives,
  and certify the required local data at primes dividing each order.
  Separate machine-checkable finite tables from transport and coverage.
\end{enumerate}

\begin{gate}[Finite exceptional families]
Each family supplies
\path{PrimeSupportedOneUniversalCoverOneSylowFamilyVerification} with no
table or projective hypothesis.  Representative transport is tested on a
second isomorphic cover and a conjugate Sylow subgroup.
\end{gate}

\section{P3: finite reductive groups and final type \(D\)}

No directory currently implements this phase.  The order below mirrors the
actual dependency chain in Cabanes--Sp\"ath \cite{CabanesSpath2026}; theorem
numbers are traceability labels, not substitute axioms.

\subsection{P3.A: finite-reductive foundations}

\begin{enumerate}[label=\textbf{P3.A\arabic*.},leftmargin=*]
  \item Define connected reductive groups with Frobenius/Steinberg maps,
  fixed-point finite groups, regular embeddings, dual data, and transport
  between isogenous models.
  \item Develop Sylow \(d\)-tori, their normalizers and centralizers,
  relative Weyl groups, and the arithmetic \(d=d_\ell(q)\).
  \item Formalize the global equivariant character parametrization
  \(\mathrm A(\infty)\) (Theorem~2.18) explicitly.  Do not hide it inside
  the local criterion.
  \item Define \(\mathrm A(d)\) and \(\mathrm B(d)\), their transversals,
  extension maps, maximal extendibility, and equivariance actions.
\end{enumerate}

\begin{gate}[P3.A]
The local criterion corresponding to Theorem~2.21 is a proved theorem with
the exact premise \(\ell\nmid 6q\), and \(\mathrm A(\infty)\),
\(\mathrm A(d)\), and \(\mathrm B(d)\) appear as separately proved inputs.
\end{gate}

\subsection{P3.B: family cases before the final type \(D\) seam}

\begin{enumerate}[label=\textbf{P3.B\arabic*.},leftmargin=*]
  \item Encode defining-characteristic, small-prime, exceptional, and
  classical-family results with their exact hypotheses.
  \item Record Theorem~2.11 accurately: for \(\ell\geq5\) the remaining
  nondefining-characteristic exclusions are type \(D\) and
  \({}^{2}D\); do not state it as all Lie types.
  \item Instantiate the one-cover/one-Sylow prime-supported family
  interface for every completed family.
\end{enumerate}

\begin{gate}[P3.B]
All Lie-type representatives not delegated to the final \(D/{}^2D\)
module have local data, and the union theorem exposes exactly the remaining
type-\(D\) predicate.
\end{gate}

\subsection{P3.C: final type \(D\)}

\begin{enumerate}[label=\textbf{P3.C\arabic*.},leftmargin=*]
  \item Formalize the doubly regular and rank-four cases (Theorem~4.1).
  \item Formalize the arithmetic trichotomy of Lemma~5.8: doubly regular,
  cyclotomic multiplicity at most one, or the disconnected maximal-rank
  subgroup with two type-\(D\) factors.
  \item Construct the subgroup \(M\), its normalizers and Sylow
  \(d\)-tori, and the stable-transversal result through Theorem~5.20.
  \item Formalize the Section~6 extension maps and maximal-extendibility
  results, culminating in Theorem~6.9 and the local conditions.
  \item Assemble Theorem~6.12 with its exact hypothesis
  \(\ell\nmid 2q\), including rank four, \(d=1,2\), cyclic-Sylow, doubly
  regular, and residual cases.
\end{enumerate}

\begin{gate}[P3.C]
A semantic theorem such as \path{TypeD.inductiveMcKay} covers exactly
Theorem~6.12 without assumed local conditions, extension maps, or
classification facts.  It then instantiates the prime-supported Lie-family
interface.
\end{gate}

\section{P4: classification coverage and universal-cover assembly}

\subsection*{Low-level plan}

\begin{enumerate}[label=\textbf{P4.\arabic*.},leftmargin=*]
  \item Define concrete, isomorphism-invariant finite-Lie-type and sporadic
  families beside the existing alternating family.
  \item Formalize a CFSG representative exhaustion or import a compatible,
  licensed, kernel-checked one.  The exact endpoint is an inhabitant of
  \path{CFSGFamilyCoverage}, not a comment citing CFSG.
  \item Use
  \path{PrimeSupportedCFSGFamilyVerification}: verify only primes dividing
  the representative simple-group order.
  \item Use the existing transport boundary correctly: one representative
  per simple-group isomorphism class, one universal cover of that
  representative, and one Sylow subgroup suffice.  Transport to arbitrary
  targets, covers, and Sylow subgroups only after the representative theorem
  is complete.
  \item Assemble
  \path{UniversalCoverInductiveMcKayHypothesis}; use its existing descent
  to \path{QuasisimpleInductiveMcKayHypothesis} for the current final
  cardinality reduction.
\end{enumerate}

\begin{gate}[P4]
The quasisimple or universal-cover hypothesis is produced by a closed Lean
term.  \texttt{\#print axioms} shows no project-specific classification or
family assumption, and every simple-family disjunct is discharged.
\end{gate}

\section{P5: universe transport, final assembly, and audit}

\subsection*{Low-level plan}

\begin{enumerate}[label=\textbf{P5.\arabic*.},leftmargin=*]
  \item Compose the closed P1 counting theorem and P4 quasisimple theorem
  through
  \path{statement_of_ordinaryInertiaFibreCardinality_quasisimple}.
  \item The conditional final wrappers already lift their universe-zero
  composites with \path{Statement.of_universeZero} from
  \path{Proof/SmallModel.lean}.  Add a regression theorem showing that a
  closed wrapper works for a genuinely higher-universe finite group.
  \item Export \path{McKayConjecture.mckay} at the exact type in
  Section~1.
  \item Run the full build, placeholder and heartbeat audits, dependency
  pin check, and \texttt{\#print axioms McKayConjecture.mckay}.
  \item Build both TeX documents and update this status matrix from
  declarations, not intentions.
\end{enumerate}

\begin{gate}[P5: completion]
All of the following succeed from a fresh checkout at the pinned hashes:
\begin{lstlisting}
lake exe cache get
lake build
rg -n '^[[:space:]]*(axiom|sorry|admit)|maxHeartbeats' \
  McKayConjecture
#check McKayConjecture.mckay
#print axioms McKayConjecture.mckay
latexmk -pdf -interaction=nonstopmode -halt-on-error \
  docs/mckay_proof.tex
latexmk -pdf -interaction=nonstopmode -halt-on-error \
  docs/formalization_blueprint.tex
\end{lstlisting}
The theorem has no extra hypothesis, the full build is green, and the
axiom report contains no project-specific mathematical assumption.
\end{gate}

\section{File organization and milestone naming}

\subsection{Critical existing files}

\begin{description}[style=nextline,leftmargin=0pt]
  \item[\path{Character/ActionInvariantOver.lean}]
  Ordinary and \(p'\)-degree invariant lying-over types, self-conjugation
  erasure, and the conditional degree-erasure adapter.
  \item[\path{InductiveMcKay/NormalSubgroupCentralIntersectionCardinalityReduction.lean}]
  Exact cardinality hypothesis and conversion to fibre equivalences.
  \item[\path{InductiveMcKay/NormalSubgroupCentralIntersectionOrdinaryCardinalityReduction.lean}]
  Separate ordinary-count and automatic-\(q'\)-degree hypotheses, plus
  their checked composition into the prime-to-\(q\) boundary.
  \item[\path{InductiveMcKay/FinalCardinalityReduction.lean}]
  Universe-polymorphic conditional composition with quasisimple or
  universal-cover data, in both prime-to-\(q\) and
  ordinary-count-plus-degree forms.
  \item[\path{InductiveMcKay/NormalSubgroupCentralIntersectionJoinCenterFibreReduction.lean}]
  Why arbitrary finite fibre equivalences suffice after adjoining the
  ambient center.
  \item[\path{Character/CliffordPPrimeCorrespondence.lean}]
  Clifford equivalence on the \(p'\)-degree inertia fibres.
  \item[\path{InductiveMcKay/UniversalCoverFamilyReduction.lean}]
  One-cover/one-Sylow transport and CFSG-family interfaces.
\end{description}

\subsection{Planned leaf files}

\begin{description}[style=nextline,leftmargin=0pt]
  \item[\path{Character/GoodElement.lean}]
  Good elements/classes and transport.
  \item[\path{Character/GallagherInvariantCount.lean}]
  Theorem~1.3 and its inertia version, using the fixed-character subtypes
  from \path{Character/ActionInvariantOver.lean}.
  \item[\path{Character/OkuyamaWajimaInvariantCount.lean}]
  Theorem~2.1, Corollary~2.3, Lemma~2.2, and per-parameter Equation (2.1).
  \item[\path{InductiveMcKay/NormalSubgroupCentralIntersectionFixedCount.lean}]
  Specialization to the exact existing cardinality boundary.
  \item[\path{FiniteReductive/...}, \path{TypeD/...}, \path{CFSG/...}]
  Add only when their first independently compiling semantic object exists;
  do not create placeholder directory trees.
\end{description}

Use semantic theorem names for stable APIs.  Paper-number names may be
aliases next to the proof.  Whenever an intermediate result can be stated
using current mathlib or repository types, state it immediately: this makes
the dependency graph executable, gives Lean's elaborator a small target,
and prevents a long proof from hiding a mismatch at the final boundary.

\section{Risk register}

\begin{longtable}{p{0.24\textwidth}p{0.34\textwidth}p{0.33\textwidth}}
\toprule
Risk & Failure mode & Mitigation/gate\\
\midrule
Aggregate/per-fibre confusion &
Theorem~2.4 total count is used to claim each inertia-fibre equality. &
Export and test Equation (2.1) before summing; P1.D is mandatory.\\

Ordinary/\(p'\)-degree confusion &
Equation (2.1) is treated as an equality of \(p'\)-degree fibres even
though it counts ordinary invariant characters. &
Keep both subtypes, prove the automatic-degree theorem on each side, and
cross the conditional degree-erasure equivalence before Clifford
transport.\\

Wrong character on the target &
The abstract Glauberman--Gallagher character is only propositionally or
orbitwise related to \path{internalTheorem44CharacterEquiv}. &
Transport the quotient parameter with
\path{productLinearQuotientEquiv}; prove the literal character bridge and
its covariance in P1.E1.\\

Universe-zero endpoint &
The final theorem proves \path{Statement} only for \texttt{Type}. &
Use the existing \path{Statement.of_universeZero} theorem in the exported
wrapper and compile a genuinely higher-universe specialization.\\

\(A_6\) numerical/full confusion &
The completed 31-row table is reported as full inductive McKay. &
Require all three full-stabilizer propositions and rowwise projective
completions before constructing
\path{AlternatingSixExceptionalPrimeLocalData}.\\

Cohomology/literal-factor confusion &
Equal \(H^2\)-classes are supplied where a literal factor equality is
required. &
Use a checked gauge/rescaling and finish the exact
\path{factorAgreement} field.\\

CFSG hidden as an interface &
\path{CFSGFamilyCoverage} remains a theorem argument. &
P4 is complete only when the predicate has a closed inhabitant.\\

Literature theorem overstated &
Theorem~2.11 or 2.21 is applied outside its hypotheses. &
Encode exact prime and family exclusions; isolate each case in its own
compile gate.\\

Build regression &
Large generated computations or concurrent edits break an umbrella import. &
P0 preempts all mathematics; use one full build and leaf-level files.\\
\bottomrule
\end{longtable}

\section{Definition of genuine completion}

The project is complete only when the unconditional, universe-polymorphic
\path{McKayConjecture.mckay} builds from a fresh pinned checkout and its
dependency report contains no project-specific assumption.  The following
are valuable but are not completion:

\begin{itemize}
  \item the current cardinality and quasisimple conditional theorems;
  \item the numerical \(A_6\) McKay equalities;
  \item a proof of Maltempo--Vallejo Theorem~2.4 without per-\(\mu\)
  Equation (2.1);
  \item a family verification with finite table or projective data as an
  argument;
  \item a \path{CFSGFamilyCoverage} parameter; or
  \item a universe-zero theorem presented as the public
  universe-polymorphic statement.
\end{itemize}

The route in this blueprint minimizes the ordinary-group reduction burden,
but the classification and family verification remain a large formalization
project.  Keeping those boundaries explicit is part of the proof.

\begingroup
\footnotesize
\bibliographystyle{alpha}
\bibliography{references}
\endgroup

\end{document}
